\section*{Hoofdstuk \ref{ch:b3:structural-biology} --- Structuurbiologie van eiwitten}

\begin{solution}{exo:b3:structural-biology:1}
Primair: de volgorde van de $141$ residuen van een $\alpha$-keten.
Secundair: haar acht $\alpha$-helices. Tertiair: de globinevouwing van
één keten om haar heem gewikkeld. Quaternair: het
$\alpha_{2}\beta_{2}$-tetrameer en de manier waarop zijn dimeren ten
opzichte van elkaar draaien.
\end{solution}

\begin{solution}{exo:b3:structural-biology:2}
$36\times \qty{0.15}{nm} = \qty{5.4}{nm}$; $36/3.6 = 10$ windingen;
$36 - 4 = 32$ waterstofbruggen.
\end{solution}

\begin{solution}{exo:b3:structural-biology:3}
Dat de natieve structuur, met de juiste paring van vier zwavelbruggen
uit $105$ mogelijkheden, spontaan uit de ongevouwen keten wordt
teruggevonden: de sequentie bevat alle benodigde informatie, en de
natieve toestand is het thermodynamische minimum. Het in een
reageerbuis met zuiver eiwit doen sloot elke matrijs, elk enzym en elke
celmachinerie als bron van die informatie uit.
\end{solution}

\begin{solution}{exo:b3:structural-biology:4}
NMR: kleine eiwitten, onder ongeveer \qty{40}{kDa}, in oplossing.
Kristallografie: elke grootte die kristalliseert, van peptiden tot
ribosomen. \omterm{def:b3:structural-biology:methods}{Cryo-elektronenmicroscopie}: grote complexen, boven ongeveer
\qty{50}{kDa}, zonder kristallen. NMR meldt beweging rechtstreeks (en
cryo-EM kan moleculen in conformatieklassen indelen); een
kristalstructuur is een statisch gemiddelde.
\end{solution}

\begin{solution}{exo:b3:structural-biology:5}
$2^{150} \approx 1.4\times 10^{45}$ conformaties, $1.4\times 10^{33}$
s $\approx 5\times 10^{25}$ jaar. $3^{20} \approx 3.5\times 10^{9}$,
$3.5$ ms: het peptide zou binnen een seconde uitputtend kunnen zoeken;
het eiwit niet in de ouderdom van het heelal.
\end{solution}

\begin{solution}{exo:b3:structural-biology:6}
$Y = [\alpha(1+\alpha) + Lc\alpha(1+c\alpha)]/[(1+\alpha)^{2} +
L(1+c\alpha)^{2}]$. $\alpha = 1$: $3.01/106 = 0.028$ (één plaats
$0.50$). $\alpha = 10$: $121/242 = 0.50$ ($0.91$). $\alpha = 100$:
$\num{10300}/\num{10601} = 0.97$ ($0.99$). De MWC-curve blijft bij
weinig ligand ver achter bij die van één plaats en stijgt dan steil ---
van $0.03$ naar $0.97$ over twee decaden waar de ene plaats van $0.5$
naar $0.99$ gaat: die steilheid is de \omterm{def:b3:structural-biology:allostery}{coöperativiteit}.
\end{solution}

\begin{solution}{exo:b3:structural-biology:7}
Bisfosfoglyceraat bindt alleen de \omterm{def:b3:structural-biology:allostery}{T-toestand} (in de centrale holte) en
stabiliseert T dus: in het model verhoogt het $L$, de verhouding $T/R$
van het eiwit zonder ligand, en verschuift het de curve naar rechts ---
lagere affiniteit, meer zuurstof afgegeven bij weefseldrukken. Bewaarde
rode cellen verliezen de stof, $L$ daalt, de curve verschuift naar
links, en getransfundeerd bloed houdt zijn zuurstof vast in plaats van
haar in de weefsels af te geven, tot de cellen haar in een dag opnieuw
aanmaken.
\end{solution}

\begin{solution}{exo:b3:structural-biology:8}
$d = \lambda/(2\sin\theta) = 0.1/(2\times 0.259) = \qty{0.19}{nm}$: de
ruggengraat en elke zijketen liggen vast, afzonderlijke atomen zijn
opgelost. Voor \qty{0.12}{nm}: $\sin\theta = 0.1/0.24 = 0.417$, $\theta =
\qty{25}{\degree}$.
\end{solution}

\begin{solution}{exo:b3:structural-biology:9}
Een geladen lysine begraven tussen apolaire zijketens kost tientallen
kilojoules per mol (haar lading niet gesolvateerd, de hydrofobe
begraving van de leucine verloren), meer dan de hele $\Delta G$ van het
vouwen: het eiwit ontvouwt zich of de kern schikt zich opnieuw. Aan het
oppervlak is de lysine gehydrateerd zoals zij in de ongevouwen toestand
zou zijn en gaat er niets verloren. $\qty{-30}{kJ/mol} \approx -12RT$
betekent $K = e^{12} \approx 1.6\times
10^{5}$: op elk moment is één molecuul op \num{160000} ongevouwen, en de
hele marge staat gelijk aan twee of drie waterstofbruggen van de
honderden --- het equivalent van één mutatie.
\end{solution}

\begin{solution}{exo:b3:structural-biology:10}
De $R$-soorten tellen samen $[R_{0}](1+\alpha)^{n}$ en de $T$-soorten
$[R_{0}]L(1+c\alpha)^{n}$, dus $\bar R = (1+\alpha)^{n}/[(1+\alpha)^{n}
+ L(1+c\alpha)^{n}]$. Voor $c \to 0$ is de $T$-term $L$, en is $\bar R =
1/2$ wanneer $(1+\alpha)^{n} = L$, dat wil zeggen $\alpha = L^{1/n} - 1$.
Voor hemoglobine is $(3\times 10^{5})^{1/4} = 23.4$, dus $\alpha \approx 22$:
\qty{22}{mmHg}, dicht bij de $P_{50}$ van \qty{26}{mmHg} --- de
conformatieomschakeling en de halfverzadiging vallen bijna samen.
\end{solution}

\begin{solution}{exo:b3:structural-biology:11}
Zij stuitte op weerstand omdat elke bekende besmettelijke verwekker
zich via nucleïnezuur vermenigvuldigde, en een eiwit zijn sequentie
niet kan kopiëren. Een nucleïnezuur waarvan zou blijken dat het nodig
is, of gezuiverd eiwit dat de ziekte niet zou kunnen overdragen, zou
haar weerleggen. Steun: de besmettelijkheid overleeft nucleasen,
ultraviolette en ioniserende straling bij doses die elk genoom
vernietigen, maar wordt vernietigd door eiwitdenaturerende middelen en
proteasen; recombinant PrP dat in vitro tot vezels is gevouwen draagt
de ziekte over op dieren, met dezelfde stameigenschappen bij
doorgave. Een kiem is nodig omdat de omzetting van één normaal monomeer
alleen onbetaalbaar traag is --- de kern van de
cross-$\beta$-structuur moet eerst ontstaan, en dat is de
snelheidsbepalende gebeurtenis --- terwijl een reeds gevormd
vezeluiteinde monomeren snel omzet: de vorm wordt door gestuurde groei
gekopieerd.
\end{solution}

\begin{solution}{exo:b3:structural-biology:12}
Vele sequenties vouwen zich tot dezelfde structuur: wat een \omterm{def:b3:structural-biology:levels}{vouwing}
nodig heeft is een patroon van hydrofobe en polaire posities, enkele
glycines en prolines bij de bochten, en de residuen die het werk doen;
al het overige aan het oppervlak kan zonder straf veranderen, zodat
mutaties zich daar ophopen terwijl de selectie de \omterm{def:b3:structural-biology:levels}{vouwing} en de functie
behoudt. De bewaarde residuen zijn dus de kern (eerder als patroon dan
als identiteiten), de katalytische en bindende residuen, en de
structurele glycines, prolines en cysteïnes. Profielmethoden wegen juist
die kolommen en herkennen zo familieleden bij \qty{15}{\%} identiteit;
het ontwerpen van eiwitten draait de redenering om, kiest een hydrofoob
patroon voor een beoogde \omterm{def:b3:structural-biology:levels}{vouwing} en laat het oppervlak vrij.
\end{solution}

\begin{solution}{pb:b3:structural-biology:1}
\textbf{1.} $0.75\times 150 \approx 112$ residuen in helices; $112\times
\qty{0.15}{nm} = \qty{17}{nm}$; $112/3.6 \approx 31$ windingen.
\textbf{2.} $112 - 8\times 4 = 80$ waterstofbruggen.
\textbf{3.} Massa $150\times 110 = \qty{16500}{Da} = 2.7\times 10^{-20}$
g; volume $2.7\times 10^{-20}\times 0.73 = 2.0\times 10^{-20}$ cm$^{3}$
$= \qty{20}{nm^3}$; straal $(3V/4\pi)^{1/3} = \qty{1.7}{nm}$.
\textbf{4.} Gestrekt $150\times 0.36 = \qty{54}{nm}$ tegen een diameter
van \qty{3.4}{nm}: een factor $16$.
\textbf{5.} De residuen in helices vormen \qty{75}{\%} van de gestrekte
lengte, $112\times 0.36 = \qty{40}{nm}$; de helix perst ze samen tot
\qty{17}{nm}, een factor $0.36/0.15 = 2.4$.
\textbf{6.} $40\times 4 = \qty{160}{kJ/mol}$ gewonnen, $150\times 0.9 =
\qty{135}{kJ/mol}$ verloren: netto $\Delta G \approx \qty{-25}{kJ/mol}$,
in het marginale bereik.
\textbf{7.} $3^{150} \approx 4\times 10^{71}$; $4\times 10^{59}$ s
$\approx 10^{52}$ jaar.
\textbf{8.} $10^{-2}\times 10^{12} = 10^{10}$ conformaties, een fractie
$3\times 10^{-62}$ van de telling.
\textbf{9.} Helices in \qty{1}{\micro s} (parallel); $8! = \num{40320}$
pakkingen bij $10^{6}$ per seconde kosten \qty{40}{ms}: samen ongeveer
\qty{40}{ms}, de waargenomen orde van grootte. Ja: de delen
afzonderlijk oplossen en daarna het geheel is een trechter --- de
zoekruimte wordt ontbonden, niet uitgeput.
\textbf{10.} Verwacht aantal cycli $1/0.3 = 3.3$, ongeveer \qty{33}{s};
na zes cycli is $0.7^{6} = 0.12$ nog ongevouwen.
\textbf{11.} $\Delta G = -RT\ln K$ met $K = 10^{4} \to 10^{2}$: van
$-23.7$ naar \qty{-11.9}{kJ/mol}, een verlies van \qty{12}{kJ/mol}
stabiliteit ($RT = \qty{2.58}{kJ/mol}$, $\ln 100 = 4.6$).
\textbf{12.} Klontering begint met een kiem, waarvan de vormingssnelheid
als een hoge macht van de concentratie van de klonteringsgevoelige
ongevouwen soort stijgt; honderd keer meer van die soort verhoogt de
kans dat er binnen een mensenleven een kiem ontstaat met orden van
grootte, en zodra er een kiem is verloopt de groei snel en
zelfsturend.
\textbf{13.} $\alpha = 100$: $Y = 1.054\times 10^{8}/1.089\times 10^{8} =
0.97$. $\alpha = 20$: teller $185\,220 + 103\,680 = 288\,900$,
noemer $194\,481 + 622\,080 = 816\,561$: $Y = 0.35$.
\textbf{14.} $\bar R = 1.041\times 10^{8}/1.089\times 10^{8} = 0.96$ in
de long; $194\,481/816\,561 = 0.24$ in de spier.
\textbf{15.} $0.97 - 0.35 = 0.62$ van de capaciteit gelost (twee derde
van de lading). Eén plaats: $100/126 - 20/46 = 0.79 - 0.43 = 0.36$.
\textbf{16.} Capaciteit $150\times 1.34 = \qty{200}{mL/L}$; naar de
spier $0.62\times 200 \approx \qty{125}{mL/L}$. In rust $(0.97 - 0.78)\times
200 = \qty{38}{mL/L}$; $250/38 \approx \qty{6.5}{L/min}$
bloedstroom --- de orde van een hartminuutvolume in rust.
\textbf{17.} $L = 3\times 10^{6}$, $\alpha = 40$: teller $2.76\times
10^{6} + 3.29\times 10^{6} = 6.05\times 10^{6}$, noemer $2.83\times
10^{6} + 11.5\times 10^{6} = 14.4\times 10^{6}$: $Y = 0.42$ in plaats
van $0.78$ --- veel meer zuurstof gelost in rust, de aanpassing van
bloed aan hoogte en van spier aan arbeid.
\textbf{18.} Foetaal, $L = 3\times 10^{4}$, $\alpha = 30$: teller
$893\,730 + 19\,770 = 913\,500$, noemer $923\,520 + 85\,680 =
1\,009\,200$: $Y = 0.91$. Bij de moeder op dezelfde druk: teller
$893\,730 + 197\,730 = 1\,091\,460$, noemer $923\,520 + 856\,830
= 1\,780\,350$: $Y = 0.61$. Bij gelijke druk houdt het foetale bloed
meer zuurstof vast, zodat er zuurstof van moeder naar foetus over de
placenta stroomt.
\textbf{19.} $(10^{5}\,\text{nm}/6\,\text{nm})^{3} = 4.6\times 10^{12}$
cellen; $1.9\times 10^{13}$ moleculen.
\textbf{20.} $d = 0.1/(2\sin 14.5^{\circ}) = \qty{0.20}{nm}$: ja, elke
zijketen en elk atoom liggen vast.
\textbf{21.} $\sin\theta = 0.1/0.3$, $\theta = \qty{19.5}{\degree}$. In
een wanordelijk kristal staan de moleculen niet precies in het gelid,
tellen de golven die onder een grote hoek worden verstrooid --- en die
de fijne details bevatten --- niet meer in fase op, en vervagen de
vlekken bij grote hoeken; alleen de vlekken met lage \omterm{def:b3:structural-biology:methods}{resolutie} blijven
over.
\textbf{22.} $d$ halveren vergt tweemaal het signaal, dus viermaal
zoveel deeltjes: ongeveer \num{40000} (in de praktijk meer, omdat er
andere fouten bij komen).
\textbf{23.} Membraaneiwitten, die zelden uit detergensoplossingen
kristalliseren en in een lipidennanoschijf kunnen worden afgebeeld;
grote buigzame complexen --- ribosomen, spliceosomen, ionkanalen met
hun regulatoren --- waarvan de heterogeniteit de kristallisatie
verhindert maar waarvan de deeltjes in conformatieklassen kunnen worden
gesorteerd.
\textbf{24.} De gebonden heem en zuurstof en hun precieze meetkunde,
het geordende water en de ionen, de tweede conformatie (T tegenover R)
en het bewijs dat de \omterm{def:b3:structural-biology:levels}{vouwing} klopt --- een voorspelling is één enkel
model zonder ligand en zonder experimentele toets.
\textbf{25.} Lossen van $0.62$ van de capaciteit (één plaats $0.36$); de
tijd van Levinthal ongeveer $10^{52}$ jaar; een kaart bij
\qty{0.20}{nm}.
\end{solution}
